% 汇盈探店 · 管理员后台使用手册（xelatex 编译）
\documentclass[12pt,a4paper,fontset=none]{ctexart}
\newcommand{\manualtitle}{管理员后台使用手册}
% 共享排版样式：管理员手册 / 探店员手册 通用前言（xelatex + ctex 编译）
\usepackage[margin=2.2cm,headheight=14pt]{geometry}
\usepackage{xcolor}
\usepackage{hyperref}
\usepackage{booktabs}
\usepackage{tabularx}
\usepackage{array}
\usepackage{enumitem}
\usepackage{fancyhdr}
\usepackage{needspace}
\usepackage{colortbl}

\hypersetup{colorlinks=true,linkcolor=black,urlcolor=blue,pdftitle=\manualtitle}

% 字体（macOS 中文字体）
\setCJKmainfont{Songti SC}[BoldFont={Songti SC Bold}]
\setCJKsansfont{Hiragino Sans GB}[BoldFont={Hiragino Sans GB W6}]
\setCJKmonofont{Menlo}
\setmainfont{Times New Roman}

% 品牌色
\definecolor{brand}{RGB}{242,95,66}
\definecolor{brandlight}{RGB}{255,240,236}
\definecolor{graytext}{RGB}{138,141,146}

% 标题样式：黑体 + 品牌色
\ctexset{
  section/format={\Large\bfseries\sffamily\color{brand}},
  subsection/format={\large\bfseries\sffamily\color{black}},
  subsubsection/format={\normalsize\bfseries\sffamily},
}

% 页眉页脚
\pagestyle{fancy}
\fancyhf{}
\fancyhead[L]{\small\sffamily\color{graytext}\manualtitle}
\fancyhead[R]{\small\sffamily\color{graytext}\thepage}
\renewcommand{\headrulewidth}{0.4pt}
\renewcommand{\headrule}{\color{brand!60}\hrule width\headwidth height 0.4pt}

% 提示框：\tipbox{标题}{内容}
\newcommand{\tipbox}[2]{%
  \par\vspace{0.5em}\noindent
  \fcolorbox{brand}{brandlight}{\begin{minipage}{\dimexpr\linewidth-2\fboxsep-2\fboxrule\relax}
    \textbf{\color{brand}#1}\par\smallskip #2
  \end{minipage}}\par\vspace{0.5em}}

% 警示框（红色）：\warnbox{内容}
\newcommand{\warnbox}[1]{%
  \par\vspace{0.5em}\noindent
  \fcolorbox{red!70}{red!8}{\begin{minipage}{\dimexpr\linewidth-2\fboxsep-2\fboxrule\relax}
    \textbf{\color{red!70!black}注意}\par\smallskip #1
  \end{minipage}}\par\vspace{0.5em}}

% 步骤列表：紧凑
\setlist[enumerate]{itemsep=2pt,topsep=4pt,leftmargin=2.2em}
\setlist[itemize]{itemsep=2pt,topsep=4pt,leftmargin=1.8em}

% 表格默认样式
\renewcommand{\arraystretch}{1.25}

% 封面
\newcommand{\manualcover}[4]{%
  \begin{titlepage}
    \centering
    \vspace*{2.5cm}
    {\color{brand}\rule{6cm}{2.5pt}}\\[1.2cm]
    {\Huge\bfseries\sffamily #1\par}
    \vspace{0.8cm}
    {\LARGE\sffamily #2\par}
    \vspace{2.5cm}
    {\large\color{graytext} 汇盈探店 · 门店服务品质暗访系统\par}
    \vspace{0.4cm}
    {\large\color{graytext} #3\par}
    \vfill
    {\color{brand}\rule{6cm}{2.5pt}}\\[0.6cm]
    {\normalsize\color{graytext} #4\par}
  \end{titlepage}}


\begin{document}

\manualcover{管理员后台使用手册}{汇盈探店 · 运营管理中心}{版本 1.0 · 2026 年 8 月}{本手册面向运营管理与复核人员，覆盖登录、门店、题库、探店任务、用户、提交记录、导出与审计的全部日常操作。}

\tableofcontents
\newpage

\section{手册说明}

\subsection{适用对象}
本手册面向「汇盈探店」运营管理中心的使用者：管理员、运营管理员、复核人员与门店负责人。你登录后台后看到的菜单由你的角色和授权可见范围决定；本手册按完整权限编写，你的界面若少某些菜单，说明你的角色没有对应权限。

\subsection{核心术语}
\begin{itemize}
\item \textbf{探店任务}：一次门店暗访计划，绑定一家门店，包含探店时间窗口、问卷截止时间和共用说明，可指派多名探店员。
\item \textbf{指派}：探店任务下分派给\emph{某一名}探店员的执行要求（问卷版本、个人时间、菜品要求、预算、小票要求）。每行指派对应一份独立问卷，不允许换人接力。
\item \textbf{暗访 / 明察}：任务类型。暗访是默认执行方式；明察目前只能保存草稿，\emph{发布时会被系统拒绝}。
\item \textbf{得分 / 不得分}：问卷每题二选一。不得分题必须填写问题说明；若题库配置了证据要求，还需按题上传图片或音频。
\item \textbf{问题证据}：按题配置的图片（照片最多张数）与音频（录音最长秒数）要求，只在「不得分」时触发。
\item \textbf{打卡（check-in）}：探店员开始任务前必须完成的到店定位核验；无法到店时管理员可创建临时打卡豁免。
\item \textbf{可见范围}：管理员账号被授权的可见范围（全部 / 按区域 / 按门店），列表与导出都会受其约束。
\end{itemize}

\section{登录与会话}

\subsection{登录后台}
打开运营管理中心后，在登录页输入「管理员账号」与「密码」，点击「登录」。

\tipbox{密码错误的提示}{出于安全考虑，账号或密码错误时页面统一显示「登录已失效，请重新登录」，不会区分是账号错还是密码错。}

\begin{itemize}
\item 连续输错会触发限流，提示「操作过于频繁，请稍后重试」，稍等一分钟再试。
\item 请求超过 12 秒或网络不通时，分别提示「请求超时，请检查网络后重试」「暂时无法连接服务，请检查网络后重试」。
\item 管理员账号由管理方统一创建与配置，不提供自助注册。
\item \textbf{只有管理员账号可以登录后台}：探店员账号登录时会提示「该账号没有后台管理权限，请使用探店员小程序」。管理员账号同时也可以登录探店员小程序使用。
\end{itemize}

\subsection{会话有效期}
\begin{itemize}
\item \textbf{滑动续期}：只要你持续操作，每次有效请求都会把空闲窗口刷新为 \textbf{30 分钟}，日常使用不会掉线。
\item \textbf{绝对上限}：无论是否活跃，单次登录最长 \textbf{12 小时}，到期必须重新登录。
\item 后台任意操作返回「登录已失效」时会自动回到登录页，重新登录即可。
\end{itemize}

\subsection{退出登录}
点击侧边栏底部用户信息右侧的「退出登录」图标按钮。

\section{界面总览}

\subsection{侧边栏菜单}
登录后左侧为功能菜单，顶部为品牌标识，底部为你的姓名与角色。菜单与角色对应关系如下：

\begin{center}
\begin{tabular}{ll}
\toprule
菜单 & 可见角色 \\
\midrule
运营概览 & 全部角色 \\
探店任务 & 管理员、超级管理员、运营管理员 \\
用户管理 & 管理员、超级管理员、运营管理员 \\
门店管理 & 管理员、超级管理员、运营管理员、门店负责人 \\
问卷题库 & 管理员、超级管理员、运营管理员 \\
提交记录 & 全部角色 \\
导出中心 & 管理员、超级管理员、运营管理员、复核人员 \\
操作审计 & 管理员、超级管理员、运营管理员、复核人员 \\
\bottomrule
\end{tabular}
\end{center}

\warnbox{菜单隐藏只影响使用体验，不代表权限开放。所有操作仍由系统再次鉴权，越权访问会提示「你没有权限进行此操作」。}

\subsection{运营概览（数据看板）}
登录后默认进入运营概览页，页面按你的角色和可见范围实时统计并展示数据卡片，点击卡片直接进入对应功能：

\begin{itemize}
\item \textbf{待审核申请}：待审批的探店员注册申请（进入用户管理）。
\item \textbf{待发布任务}：尚未发布的探店任务（进入探店任务）。
\item \textbf{进行中任务}：正在执行的任务数。
\item \textbf{逾期指派}：已逾期的探店员指派总数。
\item \textbf{已发布问卷}：已发布的问卷版本数（进入问卷题库）。
\item \textbf{门店}：门店总数。
\item \textbf{累计提交}：探店员已提交的问卷总数（进入提交记录）。
\end{itemize}

下方为「最近提交」面板：按提交时间倒序展示最近 5 份完成的问卷（探店员、任务、得分率、提交时间），点击「查看全部」进入提交记录；尚无提交时显示空态提示。某项数据因权限不可见时，对应卡片显示「—」，不影响其他卡片。

\subsection{通用交互}
\begin{itemize}
\item \textbf{确认框}：执行发布、取消、停用等操作前会弹出确认对话框，需点「确认」才真正执行；危险操作（取消、删除、停用、不批准）的确认按钮为红色。提交中按钮显示「处理中…」，此时无法取消。
\item \textbf{提示条（Toast）}：操作结果在屏幕底部居中提示——成功为绿色对勾，失败为红色警告图标；提示不会自动消失，可点右侧「关闭提示」或等它被下一条覆盖。
\item \textbf{状态徽章}：列表中的状态以彩色徽章显示，中文含义见附录 A。
\item \textbf{时间格式}：后台统一显示为本地时间，形如「2026/8/15 14:30:00」，空值显示「—」。
\item \textbf{列表空态}：无数据显示「当前没有××」，加载失败时提供「重新加载」按钮。
\end{itemize}

\section{门店管理}

门店是探店任务的落点：创建任务前，必须先建好门店并完成位置校准。

\subsection{门店列表}
列表展示：门店（名称与编号）、区域、状态（已启用/已停用）、打卡半径、操作。行内按钮为「编辑范围」与「删除」。

\subsection{新建 / 编辑门店}
点击「新建门店」或某门店的「编辑范围」，填写以下内容：

\begin{center}
\begin{tabular}{p{3.6cm}p{9.6cm}}
\toprule
字段 & 说明 \\
\midrule
门店编号 & 必填，1–64 字符，全局唯一 \\
门店名称 & 必填，1–200 字符 \\
省份 / 城市 / 区县 & 三级联动选择；门店必须归属到\emph{区县}级 \\
地址 & 必填，1–500 字符 \\
打卡半径（米） & 打卡判定半径，100–1000 米，默认 300 \\
允许定位误差（米） & 允许的定位误差，0–200 米，默认 200 \\
门店启用，可创建新任务 & 勾选后门店可被新任务选用；不勾选即停用 \\
\bottomrule
\end{tabular}
\end{center}

\subsection{删除门店}
门店列表每行都有「删除」（红字）：

\begin{itemize}
\item \textbf{未被任何探店任务引用}的门店可直接删除（确认框再次确认后生效）。
\item \textbf{已被任务引用}的门店删除会被拒绝，提示「该门店已被探店任务引用，无法删除，可改为停用」——此时请取消勾选「门店启用」改为停用，历史数据完整保留。
\end{itemize}

\subsection{地图选点与位置确认}
\begin{enumerate}
\item 新建门店时，地图默认定位在\textbf{昆明}（业务门店主要分布区域）；编辑已有门店时定位到该店。
\item 在「搜索门店或商场」输入关键词，点击「搜索」。搜索为\textbf{全国模糊搜索}，不受当前地图范围限制：输入「建新园」会列出全国所有相关门店，输入「建新园胜利堂」直接出对应结果；点击结果后地图自动跳转到该地点并回填地址、区域和位置。
\item 也可以直接拖动地图，点击「使用地图中心位置」以地图中心作为门店位置。
\item 选点后可继续拖动地图微调，圆环半径随「打卡半径」实时变化。
\item 点击「保存门店」。

\tipbox{位置确认规则}{门店保存时带有有效位置即可正常发布任务；位置未确认的门店发布任务时会被系统提示。}
\end{enumerate}

地图相关的几种提示及处理：

\begin{itemize}
\item 「地图服务尚未配置，请联系管理员」——环境未注入腾讯地图 Key，需运维配置后刷新。
\item 「地图加载失败，请检查网络后重试」——地图脚本加载失败，检查网络后重试。
\item 「没有找到匹配地点，请调整关键词」——腾讯地点库中无此门店；可换关键词，或拖动地图用中心点校准。
\end{itemize}

\warnbox{门店位置统一使用标准地图坐标系；打卡是否通过由系统最终判定，探店员定位超出「半径 + 允许的定位误差」范围时无法开始任务。}

\section{问卷题库}

题库以「版本」管理：新建草稿版本 → 编辑题目与证据要求 → 发布；已发布版本不可修改，修改必须先「复制为新草稿」。

\subsection{问卷版本列表}
列表显示版本名称、题目检查结果（「N 题完整」或「发布检查未通过」）、状态、创建时间。操作：「查看题目」；已发布版本可「复制为新草稿」；草稿版本可「完成检查并发布」或「删除」。

\tipbox{删除问卷版本}{只有\textbf{草稿}版本且未被任何指派引用时可以删除。已发布版本没有删除入口——如不需要，请复制新草稿后放弃，源版本保持不可修改。}

\subsection{新建问卷}
点击「新建问卷」，填写问卷名称（1–64 字符）。编辑器默认给出 \textbf{3 道空题}：

\begin{itemize}
\item \textbf{添加 / 删除题目}：底部「添加题目（N/50）」逐题添加，最多 50 道；每题右侧「−」删除，至少保留 1 道。删除后题号自动重新编号。
\item \textbf{类别}：环境 / 服务 / 产品。按题号自动预填（1–20 环境、21–42 服务、43 以后产品），可手动调整。
\item \textbf{题干}：必填，1–500 字符。
\item \textbf{分值}：0–10 分，默认 2 分。
\item \textbf{照片 / 录音}：勾选「照片」后出现「照片最多张数」（1–9 张，默认 3）；勾选「录音」后出现「录音最长秒数」（1–60 秒，默认 10）。
\end{itemize}

点击「创建问卷草稿」保存。\textbf{只提交题干非空的题目}，请至少填写一道题。

\subsection{题目设置与证据配置（版本详情页）}
草稿版本详情页左侧为题号列表，右侧为选中题的设置：

\begin{itemize}
\item \textbf{题目内容}：题干（草稿可改，已发布不可修改）。
\item \textbf{分值}：0–10。
\item \textbf{不得分证据要求}：
  \begin{itemize}
  \item 开关「不得分时要求上传证据」——开启后至少勾选「需要图片」或「需要音频」之一（可都选）。
  \item 需要图片：探店员不得分时必须至少上传 1 张照片（最多按「照片最多」）。
  \item 需要音频：探店员不得分时必须至少录制 1 段音频（最长按「录音最长」）。
  \end{itemize}
\end{itemize}

\tipbox{规则提醒}{「问题说明」在不得分时永远必填，与证据开关无关；证据要求只在题目标记「不得分」时生效，得分题不要求证据。录音每道题最多保留 3 段。}

每道题修改后点击「保存题目配置」。

\subsection{发布检查与发布}
草稿版本右上角「发布版本」。发布前系统逐项检查并显示通过/未通过：

\begin{enumerate}
\item \textbf{启用题数}：1–50 道均可发布。
\item \textbf{题号连续}：题号必须从 1 开始连续编号（删除题目后会自动重排）。
\item \textbf{证据配置}：开启「要求上传证据」的题，必须勾选了图片或音频之一。
\end{enumerate}

三项全部通过后「发布版本」才可点击；发布后版本转为不可修改，任务指派只能选用已发布版本。

\section{探店任务}

\subsection{探店任务列表}
列表展示任务名称与编号、门店、执行窗口、指派人数、逾期数、状态。操作按钮：

\begin{itemize}
\item \textbf{查看指派}：进入任务详情。
\item \textbf{编辑}（仅待发布）：修改探店任务信息。
\item \textbf{发布}（仅待发布）：将任务与全部指派发给探店员。
\item \textbf{删除}（仅待发布，红字）：删除整条任务，其下未发布的指派一并删除；确认框会再次提示。已发布或已取消的任务没有删除入口，请用「整体取消」。
\item \textbf{整体取消}（除已完成、已取消外）：需填写取消原因，红字操作。
\end{itemize}

\tipbox{删除任务}{只有「待发布」状态的任务可以删除；任务一旦发布或取消，就只能「整体取消」保留留痕，不能再物理删除。}

\subsection{新建 / 编辑探店任务}
字段如下：

\begin{center}
\begin{tabular}{p{3.6cm}p{9.6cm}}
\toprule
字段 & 说明 \\
\midrule
探店任务名称 & 必填 \\
门店 & 只能选「已启用」的门店 \\
探店日期 & 到店执行日期（单日） \\
探店时间段 & 开始时间与结束时间，结束必须晚于开始 \\
问卷最终截止 & 问卷提交的最晚时间，\emph{不能早于探店结束时间}，允许跨天 \\
共用说明 & 选填，≤2000 字符，例如保密提示与观察重点 \\
\bottomrule
\end{tabular}
\end{center}

\tipbox{两个时间的区别}{探店时间段约束「到店打卡与开始任务」；问卷最终截止约束「提交问卷」。探店结束后、截止前，探店员仍可补录提交；过了问卷截止即逾期，需要你延期后才能提交。}

\subsection{任务详情与任务指派}
任务详情页顶部为状态、门店位置、指派数、逾期数四张摘要卡；下方为「探店员指派表」。

\subsubsection{添加探店员（仅待发布）}
点击「添加探店员」，填写：

\begin{itemize}
\item \textbf{探店员}：仅列出「已启用」的探员账号。
\item \textbf{任务类型}：暗访 / 明察（明察只能保存草稿，不能发布）。
\item \textbf{已发布问卷版本}：仅列已发布的题库版本。
\item \textbf{预算上限（元）}：≥0。
\item \textbf{个人开始 / 截止时间}：必须落在探店任务执行窗口之内。
\item \textbf{个人特别说明}：选填。
\item \textbf{要求上传消费小票}：勾选后探店员至少上传 1 张小票原图才能提交。
\item \textbf{菜品要求}：逐项填写菜品名称、数量、必点、允许替代；可「添加菜品」或删除。必点菜品探店员不可跳过（售罄须填原因）。
\end{itemize}

点击「保存指派」。

\subsubsection{指派行操作（按状态出现）}
\begin{itemize}
\item \textbf{编辑要求}：待发布、待开始状态可改。
\item \textbf{延期}：待开始、进行中、已逾期可延期；需输入新的截止时间与延期原因，新时间必须晚于当前时间。逾期任务延期后自动恢复。
\item \textbf{打卡豁免}：待开始且无生效豁免时可创建，最长 30 分钟，探店员首次开始问卷时自动消费。
\item \textbf{豁免记录}：查看该探店员的历史豁免（有效/已使用/已撤销/已过期），可撤销生效中的豁免。
\item \textbf{删除}：仅待发布且尚未建立问卷时可用。
\item \textbf{取消}：除待发布、已提交、已取消外的状态可单条取消，需填原因。
\end{itemize}

\subsection{发布任务}
发布前系统会逐项检查，不满足时提示「探店任务尚未满足发布条件」并\emph{指明具体原因}，常见如下：

\begin{itemize}
\item 门店已停用 / 门店位置未校准——回「门店管理」处理。
\item 尚未添加任何探店员指派。
\item 某条指派是明察任务——发布仅支持暗访。
\item 探店员账号未启用或手机未验证。
\item 问卷版本未发布。
\item 指派时间超出探店任务执行窗口。
\item 指派数量超过上限（单探店任务最多 500 条）。
\end{itemize}

全部通过后任务发布：各指派出现在对应探店员的小程序中（第一阶段不发通知，需你线下告知探店员）。

\section{用户管理}

\subsection{探员管理与管理员管理}
用户管理页包含两个标签页：\textbf{探员管理}（含注册审批区）与 \textbf{管理员管理}。

\begin{itemize}
\item \textbf{探员管理}：顶部为「注册审批」面板（待审核 / 已批准 / 未通过 / 已停用四个状态页签），下方为探员列表（用户名、姓名、手机号、来源、状态、操作）。来源显示「注册」（小程序自助注册）或「手动创建」。
\item \textbf{管理员管理}：管理员列表，含可见范围列（全部 / 按区域 / 按门店）。
\end{itemize}

\subsection{新建账号}
点击「新建探员」或「新建管理员」：

\begin{itemize}
\item \textbf{用户名}：4–32 字符，文字/数字/下划线（仅新建）。
\item \textbf{姓名}：1–64 字符。
\item \textbf{手机号}：1 开头 11 位数字（仅新建）。
\item \textbf{账号类型}：探员 / 管理员。
\item \textbf{初始密码}：管理员设置不受复杂度限制，建议 8 位以上（编辑时为「重置密码」，留空表示不修改）。
\item \textbf{可见范围}（仅管理员）：全部 / 按区域（选到区县）/ 按门店。
\end{itemize}

\tipbox{创建即生效}{后台创建的探员账号\emph{直接激活}，无需审批即可登录小程序；管理员账号创建后即可登录后台。重置密码后，该账号的所有旧会话立即失效。}

\warnbox{可见范围建议}{「按区域 / 按门店」范围仍在完善中，受限账号目前可能被系统拒绝操作。当前版本建议新管理员统一使用「全部」范围。}

\subsection{删除账号}
探员列表与管理员列表每行都有「删除」（红字）：

\begin{itemize}
\item \textbf{仅未被任何业务数据引用}的账号可以删除（未参与过任务、审批、审计、导出等），确认框再次确认后生效；删除时其角色与可见范围记录一并清除。
\item \textbf{已参与业务数据}的账号删除会被拒绝，提示「该账号已参与业务数据，无法删除，可改为停用」。
\item 不能删除自己的账号；不能删除最后一名管理员。
\end{itemize}

\subsection{注册审批}
小程序自助注册的探员需审批后才能登录接任务：

\begin{itemize}
\item \textbf{批准}（待审核）：申请转为已启用。
\item \textbf{不批准}（待审核）：必须填写原因，申请转为未通过。
\item \textbf{重新打开}（未通过）：填原因后回到待审核，可再次审批。
\item \textbf{停用 / 启用}：停用立即撤销该账号全部会话，任务被冻结；重新启用后需重新登录。
\end{itemize}

\warnbox{小程序端当前已隐藏「注册 / 忘记密码」入口，新探员账号主要由你在后台手动创建；注册审批仅处理历史已存在的自助申请。}

\section{提交记录}

\subsection{筛选与列表}
提交记录页按「筛选条件卡片 → 提交结果 → 表格/空状态」三段式布局展示探店员已提交的问卷结果。页面内容最大宽度 1440px，超宽屏不会无限拉伸。

\subsection{筛选条件}
筛选区为独立卡片，高频字段优先：

\begin{itemize}
\item \textbf{第一排（高频）}：门店、区域、探店员、问卷状态（已提交 / 处理中 / 提交失败），均为\textbf{下拉选择}——选项来自业务数据（门店、探员、已建门店所属区域、探店任务），新建数据后自动出现，无需手动输入。
\item \textbf{提交时间}：一个时间范围（开始 — 结束），占两列宽度。
\item \textbf{更多筛选}（默认折叠）：展开后出现探店任务（下拉）、任务编号、最低得分率、最多不得分数四个低频条件。
\item \textbf{操作区}：左侧「更多筛选」开关，右侧「重置」「应用筛选」独立成行，不再挤在格子角落。
\end{itemize}

筛选为\textbf{叠加语义}：默认显示全部记录；每添加一个条件都在当前结果上继续收窄，多条件同时生效，互不覆盖。

\subsection{提交结果}
结果区上方为独立标题行：「提交结果 共 N 条」+ 右侧「导出结果」按钮（有筛选条件时显示「导出当前结果」）。表格列为提交记录（任务名与编号）、门店/区域、探店员、提交时间、得分率、不得分题数；行内「查看详情」进入详情。

\subsection{空状态}
\begin{itemize}
\item \textbf{无任何历史数据}：显示「暂无提交记录——探店员完成并提交问卷后，记录将显示在这里」。
\item \textbf{筛选导致 0 条}：显示「暂无匹配结果——尝试调整或清除当前筛选条件」，并附「清除筛选」按钮一键恢复。
\end{itemize}

\subsection{提交详情}
\begin{itemize}
\item \textbf{摘要}：问卷版本、总分、得分率、不得分题数；另有环境卫生 / 为客服务 / 产品品质三类得分率。
\item \textbf{消费小票与最终菜品比对}（有小票时显示）：小票识别结果与探店员最终填写并排展示，金额差异一目了然；缩略图可「查看短时原图」。
\item \textbf{完整题目}：50 题逐题展示得分/不得分、问题说明、证据图片与音频（可播放）。
\item \textbf{菜品记录}：菜名、状态（已点/已替代/已售罄/未点/额外菜品）、数量与评价。
\item \textbf{其他记录}：闪光点、其他建议。
\item \textbf{提交审计记录}：提交相关操作时间线。
\end{itemize}

\subsection{导出}
\begin{itemize}
\item 结果区「导出结果 / 导出当前结果」：按当前筛选条件导出（导出按钮归属结果区，与结果列表在一起）。
\item 详情页「导出这份问卷」：导出单份问卷。
\item 导出的文件按你的数据权限隐藏敏感信息，下载地址有时效（5 分钟），导出内容 7 天后自动删除；生成与下载均写入审计。
\end{itemize}

\section{导出中心}

导出中心列出你创建过的导出任务：文件、创建人、创建时间、状态。有下载地址时显示「下载」链接，否则显示「正在准备」。可在右上「创建导出」再次发起当前筛选条件的导出。

\section{操作审计}

操作审计页按时间列出业务动作记录：时间、操作人、业务动作、对象、结果。页面仅供查看，不暴露内部请求与敏感字段。

\section{常见问题与错误对照}

\begin{center}
\begin{tabular}{p{5.2cm}p{8cm}}
\toprule
提示内容 & 含义与处理 \\
\midrule
登录已失效，请重新登录 & 会话过期或账号密码错误，重新登录即可 \\
内容已更新，请刷新后重试 & 你操作期间数据被他人修改，刷新页面后重试 \\
你没有权限进行此操作 & 角色或可见范围无权限 \\
服务暂时不可用，请稍后重试 & 系统异常或数据库繁忙，稍后重试 \\
暂时无法连接服务，请检查网络后重试 & 后台连不上服务，检查网络/服务是否启动 \\
请求超时，请检查网络后重试 & 请求超过 12 秒 \\
操作过于频繁，请稍后重试 & 触发限流（登录每分钟 20 次），稍等重试 \\
该门店已被探店任务引用，无法删除，可改为停用 & 门店有任务留痕，取消勾选「门店启用」停用即可 \\
该账号已参与业务数据，无法删除，可改为停用 & 账号有任务/审批等留痕，用「停用」替代删除 \\
已发布版本不能删除 / 已发布的任务不能删除 & 发布即留痕，改用「复制为新草稿」「整体取消」 \\
探店任务尚未满足发布条件（含具体原因） & 按提示逐条处理（见 6.4 节） \\
任务已逾期，请先延期后再继续 & 需要先为对应指派延期 \\
地图服务尚未配置 & 环境未配置腾讯地图 Key，联系运维 \\
区域加载失败；请先创建区域 & 区域数据缺失，先补充区域数据 \\
操作未完成，请检查填写内容 & 表单检查不通过，按提示修正 \\
\bottomrule
\end{tabular}
\end{center}

\appendix

\section{附录 A：状态中文对照}

\begin{center}
\begin{tabular}{ll}
\toprule
状态值 & 中文显示 \\
\midrule
draft & 待发布（探店任务/问卷版本） \\
published & 待开始（探店任务）/ 已发布（问卷版本） \\
assigned & 待开始（指派） \\
in\_progress & 进行中 \\
submitted & 已提交 \\
overdue & 已逾期 \\
cancelled & 已取消 \\
pending\_approval & 待审核 \\
approved\_active & 已启用 \\
rejected & 未通过 \\
disabled & 已停用 \\
failed & 提交失败（提交筛选） \\
processing & 处理中（提交筛选） \\
\bottomrule
\end{tabular}
\end{center}

\section{附录 B：菜品与豁免状态}

\begin{center}
\begin{tabular}{ll}
\toprule
菜品状态 & 中文显示 \\
\midrule
ordered & 已点 \\
substituted & 已替代 \\
unavailable & 已售罄 \\
skipped & 未点 \\
extra & 额外菜品 \\
\midrule
豁免状态 & 中文显示 \\
\midrule
active & 有效 \\
consumed & 已使用 \\
revoked & 已撤销 \\
expired & 已过期 \\
\bottomrule
\end{tabular}
\end{center}

\end{document}
